\documentclass[12pt,a4paper]{article}

\usepackage[a4paper, margin=2.5cm]{geometry}

\usepackage[utf8]{inputenc}
\usepackage[T2A]{fontenc}
\usepackage[russian]{babel}

\usepackage{amsmath,amssymb,amsthm}
\usepackage{mathtools}
\usepackage{graphicx}
\usepackage{subcaption}
\usepackage{wrapfig}
\usepackage{booktabs}
\usepackage{caption}
\usepackage{setspace}
\usepackage{indentfirst}
\usepackage{microtype}
\usepackage{hyperref}

\DeclareMathAlphabet{\mathcal}{OMS}{cmsy}{m}{n}

\onehalfspacing
\setlength{\parindent}{1.25cm}
\setlength{\parskip}{0pt}

\captionsetup{
    font=small,
    labelfont=bf
}

\hypersetup{
    colorlinks=true,
    linkcolor=black,
    urlcolor=blue,
    citecolor=black
}

\title{Денежно-кредитная политика в условиях неполной наблюдаемости в двухактивной HANK-модели: фильтрация скрытых состояний и обучение с подкреплением в задаче выбора политики}
\author{Polina}
\date{March 2026}

\begin{document}

\maketitle
\thispagestyle{plain}

\begin{abstract}
В работе рассматривается задача денежно-кредитной политики в двухактивной HANK-модели в условиях неполной наблюдаемости состояния экономики. В отличие от стандартной постановки, в которой регулятор располагает полной информацией о состоянии системы, здесь предполагается, что истинные структурные состояния и, при необходимости, режимы экономики не наблюдаются напрямую и могут быть восстановлены лишь по шумным наблюдениям макроэкономических переменных.

Цель исследования состоит в сопоставлении двух подходов к построению денежно-кредитной политики в такой среде. Первый подход является классическим и включает представление модели в форме пространства состояний, фильтрацию скрытых состояний и последующее применение заданного правила политики. Второй подход рассматривает обучение с подкреплением как возможное расширение стандартной схемы, при котором правило выбора инструмента политики строится на основе оцененного состояния экономики.

В качестве основы используется двухактивная HANK-модель, позволяющая учитывать гетерогенность домохозяйств, различия в ликвидности активов и неоднородность реакции на изменения доходов и процентной ставки. 
Ожидаемый результат состоит в выявлении того, дает ли использование подходов на основе обучения с подкреплением преимущества по сравнению фильтрацией и фиксированным правилом политики в среде со скрытыми состояниями и ограниченной информацией.
\end{abstract}

\section{HANK-постановка как среда для анализа денежно-кредитной политики}

Опишем новую кейнсианскую модель гетерогенных агентов в духе \cite{KaplanMollViolante2018}. Эта среда позволяет одновременно учитывать номинальные жесткости, неполные рынки, незастрахованный индивидуальный риск и различия в структуре балансов домохозяйств. В отличие от стандартной новокейнсианской модели с репрезентативным агентом, такая постановка делает возможным анализ не только межвременного замещения, но и распределительных и доходных каналов денежно-кредитной политики.

Ключевая идея - домохозяйства сталкиваются с идиосинкратическими рисками дохода в условиях неполного страхования и могут сберегать в двух активах, различающихся по ликвидности и доходности: в ликвидном активе с низкой доходностью и в неликвидном активе с более высокой доходностью, использование которого связано с транзакционными издержками. В производственном секторе модели сохраняются несовершенная конкуренция и номинальные жесткости цен. Денежно-кредитная политика описывается правилом Тейлора, а государство выпускает ликвидный актив и поддерживает выполнение межвременного бюджетного ограничения за счет налогов, трансфертов и государственных расходов.
\cite{KaplanMollViolante2018}.

\subsection{Домохозяйства}

В модели рассматривается непрерывное множество домохозяйств, различающихся по объему ликвидных активов $b$, неликвидных активов $a$ и индивидуальной производительности труда $z$. В каждый момент времени агрегированное состояние домохозяйственного сектора описывается совместным распределением
\[
\mu_t(da,db,dz).
\]

Предпочтения домохозяйства задаются как
\[
\mathbb{E}_0 \int_0^\infty e^{-(\rho+\zeta)t} u(c_t,\ell_t)\,dt,
\]
где $c_t$ -- потребление, $\ell_t$ -- предложение труда, $\rho>0$ обозначает субъективную норму межвременного дисконтирования, а $\zeta>0$ -- интенсивность смертности.

Домохозяйство получает трудовой доход, доходы по активам и государственные трансферты. Ликвидный актив может использоваться как для сбережения, так и для заимствования до некоторого предела, тогда как операции с неликвидным активом сопровождаются издержками. Бюджетные ограничения имеют вид
\[
\dot b_t = (1-\tau_t) w_t z_t \ell_t + r_t^b(b_t)b_t + T_t - d_t - \chi(d_t,a_t) - c_t,
\]
\[
\dot a_t = r_t^a a_t + d_t,
\]
при ограничениях
\[
b_t \geq -\underline b, 
\qquad
a_t \geq 0.
\]

Здесь $r_t^b$ -- доходность ликвидного актива, $r_t^a$ -- доходность неликвидного актива, $d_t$ -- чистый поток средств между ликвидным и неликвидным счетом, $\chi(d_t,a_t)$ --- транзакционные издержки, $\tau_t$ --- ставка налога на трудовой доход, $T_t$ --- государственный трансферт, $w_t$ --- реальная заработная плата.

Таким образом, ликвидный актив используется прежде всего для сглаживания текущих колебаний дохода и потребления, тогда как неликвидный актив обеспечивает более высокую доходность, но связан с более ограниченной возможностью быстрого перераспределения ресурсов. Именно наличие двух активов с различной степенью ликвидности играет ключевую роль в HANK-постановке, так как позволяет модели воспроизводить реалистичное совместное распределение доходов, ликвидных и неликвидных активов, а также предельных склонностей к потреблению.

\subsection{Транзакционные издержки и различие активов}

Неликвидность актива $a$ задается через функцию транзакционных издержек
\[
\chi(d,a)=\chi_0 |d| + \chi_1 \left|\frac{d}{a}\right|^{\chi_2} a.
\]

Линейная часть издержек создает область бездействия, в которой домохозяйству невыгодно совершать мелкие перераспределения между активами. Во-вторых, выпуклая часть ограничивает скорость крупных перераспределений и не допускает мгновенных скачков портфеля. В результате доходность неликвидного актива в равновесии превышает доходность ликвидного актива, то есть
\[
r_t^a > r_t^b.
\]

Экономически это означает, что домохозяйства сталкиваются не просто с межвременным выбором между потреблением и сбережением, а с дополнительным выбором между более ликвидной и менее ликвидной формой хранения богатства. Именно это обстоятельство позволяет модели воспроизводить присутствие как домохозяйств с низким уровнем богатства и отсутствием ликвидных резервов, так и домохозяйств, располагающих значительным неликвидным богатством, но с дефицитом текущей ликвидности.

\subsection{Фирмы и номинальные жесткости}

Совокупный выпуск формируется конкурентным производителем конечного блага путем агрегирования множества промежуточных товаров. Промежуточные фирмы действуют в условиях монополии и используют капитал $k_{j,t}$ и труд $n_{j,t}$ по технологии
\[
y_{j,t}=k_{j,t}^{\alpha} n_{j,t}^{1-\alpha}.
\]

Номинальная жесткость цен задается через издержки корректировки цен по Ротембергу. Если $\pi_t=\dot P_t/P_t$ обозначает инфляцию, то новокейнсианская кривая Филлипса в непрерывном времени принимает вид
\[
\left(r_t^a-\frac{\dot Y_t}{Y_t}\right)\pi_t
=
\frac{\varepsilon}{\theta}(m_t-m^*)+\dot \pi_t,
\qquad
m^*=\frac{\varepsilon-1}{\varepsilon},
\]
где $m_t$ -- предельные издержки, $\varepsilon$ -- эластичность замещения, $\theta$ -- параметр ценовой жесткости.

Таким образом, производственный сектор модели сохраняет стандартную новокейнсианскую логику: денежно-кредитная политика влияет на реальные переменные не напрямую, а через изменение реальной процентной ставки, спроса, выпуска, занятости и предельных издержек.

\subsection{Денежно-кредитная политика и государство}

Монетарные власти управляют номинальной ставкой по правилу Тейлора
\[
i_t=\bar r^b+\phi \pi_t+\epsilon_t,
\]
где $\phi>1$, а $\epsilon_t$ интерпретируется как монетарный шок.

Реальная доходность ликвидного актива определяется через уравнение Фишера:
\[
r_t^b = i_t - \pi_t.
\]

Государство является единственным эмитентом ликвидного актива и подчиняется межвременному бюджетному ограничению
\[
\dot B_t^g + G_t + T_t
=
\tau_t \int w_t z \ell_t(a,b,z)\, d\mu_t
+
r_t^b B_t^g,
\]
где $B_t^g$ --- государственный долг, $G_t$ -- государственные расходы, $T_t$ -- трансферты, $\tau_t$ --- ставка налога на трудовой доход.

Принципиально важно, что в HANK-среде изменение процентной ставки влияет не только на межвременной выбор домохозяйств, но и на бюджет государства. Поэтому денежно-кредитная и фискальная стороны модели оказываются тесно связаны, а характер фискальной подстройки становится содержательно важной частью механизма передачи монетарного шока \cite{KaplanMollViolante2018}.

\subsection{Равновесие}

Равновесие в данной экономике задается как совокупность траекторий индивидуальных решений, цен, доходностей активов, инфляции, фискальных переменных, распределения домохозяйств и агрегированных величин, таких что:

\begin{enumerate}
    \item домохозяйства максимизируют полезность при данных ценах, налогах и трансфертах
    \item фирмы максимизируют прибыль
    \item распределение $\mu_t$ согласовано с индивидуальными оптимальными решениями
    \item государственное бюджетное ограничение выполняется
    \item все рынки очищаются
\end{enumerate}

В частности, очищение рынка ликвидного актива требует
\[
B_t^h + B_t^g = 0,
\]
где
\[
B_t^h = \int b\, d\mu_t
\]
-- совокупные holdings ликвидных активов домохозяйств.

Очищение рынка неликвидных активов записывается как
\[
K_t + q_t = A_t,
\qquad
A_t = \int a\, d\mu_t.
\]

Очищение рынка труда:
\[
N_t = \int z\, \ell_t(a,b,z)\, d\mu_t.
\]

Очищение рынка благ:
\[
Y_t = C_t + I_t + G_t + \Theta_t + \chi_t + \kappa \int \max\{-b,0\}\, d\mu_t.
\]

Тем самым агрегированная динамика определяется не только стандартными макроэкономическими ограничениями, но и эндогенной эволюцией распределения домохозяйств по состояниям.

\subsection{Монетарный шок и разложение transmission mechanism}

Далее рассматривается реакция экономики на одноразовый неожиданный экспансионистский монетарный шок. Предполагается, что экономика изначально находится в стационарном состоянии, после чего в момент $t=0$ возникает инновация в правиле денежно-кредитной политики $\epsilon_t<0$, которая затем детерминированно затухает.

Агрегированное потребление можно представить как функционал от последовательности равновесных цен и фискальных переменных:
\[
C_t(\{\Gamma_t\}_{t\ge 0}) = \int c_t(a,b,z;\{\Gamma_t\}_{t\ge 0})\, d\mu_t,
\qquad
\Gamma_t=\{r_t^b,r_t^a,w_t,\tau_t,T_t\}.
\]

Тогда полное изменение потребления в момент удара можно разложить на прямой и косвенные эффекты:
\[
dC_0
=
\int_0^\infty \frac{\partial C_0}{\partial r_t^b}dr_t^b\,dt
+
\int_0^\infty
\left(
\frac{\partial C_0}{\partial w_t}dw_t
+
\frac{\partial C_0}{\partial r_t^a}dr_t^a
+
\frac{\partial C_0}{\partial \tau_t}d\tau_t
+
\frac{\partial C_0}{\partial T_t}dT_t
\right)dt.
\]

Первое слагаемое интерпретируется как direct effect: это реакция потребления на изменение траектории доходности ликвидного актива при фиксированных заработной плате, доходности неликвидного актива и фискальной политике. Остальные слагаемые образуют indirect effects, возникающие через изменение заработной платы, доходности неликвидных активов и бюджетной реакции государства.

Именно такое разложение является центральным для HANK-анализа денежно-кредитной политики, поскольку позволяет отделить собственно межвременное замещение от общерыночных и распределительных механизмов передачи шока \cite{KaplanMollViolante2018}.

\subsection{Связь с дальнейшей задачей исследования}

Описанная HANK-постановка важна для настоящей работы как содержательная базовая среда, в которой денежно-кредитная политика зависит не только от агрегированных переменных, но и от структуры распределения домохозяйств, ликвидности их балансов и характера фискальной подстройки. Однако полная количественная реализация такой модели вычислительно сложна. Поэтому в дальнейшей части работы эта логика будет использоваться в сокращенной форме: сначала на более компактной новокейнсианской policy-среде, затем в постановке со скрытыми состояниями и, далее, при сопоставлении classical filter-plus-rule pipeline с политикой, зависящей от оценённого состояния.

\section{Базовая постановка задачи денежно-кредитной политики}

Дальнейший анализ строится вокруг задачи денежно-кредитной политики в структурной макроэкономической модели, допускающей гетерогенность агентов и наличие скрытых компонент состояния. На первом этапе, однако, полезно абстрагироваться от неполной наблюдаемости и зафиксировать базовую policy-среду при полной информации.

Пусть $x_t$ обозначает вектор состояния экономики в момент времени $t$, включающий агрегированные переменные модели и, в сокращенной форме, те характеристики распределения агентов, которые влияют на агрегированную динамику. Пусть $u_t$ обозначает инструмент денежно-кредитной политики. Тогда динамика модели может быть записана в виде
\[
x_{t+1}=f(x_t,u_t,\varepsilon_{t+1}),
\]
где $\varepsilon_{t+1}$ обозначает вектор структурных шоков.

В классической постановке орган денежно-кредитной политики наблюдает состояние $x_t$ и формирует решение по правилу
\[
u_t=\phi(x_t),
\]
где $\phi(\cdot)$ может задаваться в виде правила Тейлора или более общего параметрического правила. Качество политики оценивается с помощью стандартной функции потерь, зависящей от инфляции, разрыва выпуска и, при необходимости, дополнительных макроэкономических переменных:
\[
L_t=\pi_t^2+\lambda_x x_t^2.
\]

Именно эта постановка далее используется как исходный baseline. После этого в модель вводится информационное ограничение: policy-maker перестает наблюдать истинное состояние экономики напрямую и должен восстанавливать его по шумным наблюдениям.

\section{Переход к задаче политики при неполной наблюдаемости}

В прикладной задаче денежно-кредитной политики предположение о полном наблюдении состояния экономики является слишком сильным. На практике регулятор не наблюдает напрямую ни все структурные шоки, ни все компоненты агрегированного состояния, ни тем более полное распределение агентов. Вместо этого ему доступен лишь набор наблюдаемых макроэкономических переменных.

Пусть $x_t$ обозначает скрытое состояние экономики, а $y_t$ --- вектор наблюдаемых переменных. Тогда модель записывается в форме пространства состояний:
\[
x_{t+1}=f(x_t,u_t,\varepsilon_{t+1}), \qquad
y_t=g(x_t,\eta_t),
\]
где $\eta_t$ обозначает шум наблюдения.

В этой постановке политика должна строиться не как функция от истинного состояния $x_t$, а как функция от его оценки, полученной по истории наблюдений. Тем самым возникает классическая схема
\[
\text{структурная модель} \rightarrow \text{пространство состояний} \rightarrow \text{фильтрация} \rightarrow \text{правило политики}.
\]

Именно эта линия служит в работе базовой точкой отсчета. На ее фоне далее рассматривается вопрос о том, дает ли использование обучаемого правила политики, зависящего от оцененного состояния, преимущества по сравнению со стандартным подходом.

\section{Итоги этапа 6}

Итог этапа 6 состоит в том, что обучаемое правило полезно не как универсальная замена классическому правилу, а как более гибкий способ выбора ставки в среде со скрытыми режимами и несовершенной информацией. Оно устойчиво превосходит жестко заданное правило и особенно выигрывает там, где классическая схема ошибается в фильтрации и формировании оценки состояния, однако не дает общего преимущества над хорошо перенастроенным простым правилом.

\input{outputs/hank_regime_learning_stage6_summary/stage6_text_blocks.tex}

\end{document}
